%% ---------------------------------------------------------------------------
\section{The Right End}\label{sec:right}

The reflected triangle $R_n$ is computed from the half-line, and its right end
fires at time $2n-2$. We show that the firing cell depends only on a cone of
about $n/2$ cells of two anti-diagonals of the half-line. If two lines of
different lengths received the same values on this cone, their right ends would
fire at the same relative time, which is impossible (\cref{lem:pumping}). This
pumping lemma is a condition on the half-line alone, and it implies the
speed-$1/3$ barrier (\cref{thm:third}).

\subsection{Relative Coordinates and Determinacy}

Fix $n\ge2$ and use the relative coordinates $\tau:=t-(n-1)$ and $\kappa:=n-i$:
the relative time since the front reached the right end, and the distance from
the right end. Write $R_n(\tau,\kappa):=C_n(n-1+\tau,n-\kappa)$. Then
\[
R_n=\{(\tau,\kappa):\ 0\le\kappa\le\tau\le n-1\},
\]
the right border is at $\kappa=-1$, the left border at $\kappa=n$, the return
chain is the diagonal $\tau=\kappa$, and the firing row is $\tau=n-1$
(\cref{fig:anatomy}). The predecessors of a cell of $R_n$ lie in $R_n$, on a
border, or on one of the anti-diagonals $t+i=2n-3$ and $t+i=2n-2$, where the
line agrees with the half-line (\cref{lem:agree}). The \emph{input word} of
$R_n$ is the partial word
\[
J_n(\kappa):=\bigl(C_\infty(n-3+\kappa,\,n-\kappa),\;C_\infty(n-2+\kappa,\,n-\kappa)\bigr),
\qquad 0\le\kappa\le n-1,
\]
whose components are the \emph{lower} and \emph{upper input cells}, on the
anti-diagonals $t+i=2n-3$ and $t+i=2n-2$. (For $n=2$ and $\kappa=0$ the lower
cell lies at time $-1$; it is a predecessor of no cell of $R_n$, and we set it
to $\Lq$.)

\begin{figure}[t]
\centering
%% Figure: anatomy of the line of length n = 12, coloured by region (Lemmas 2.2, 2.4, 3.1, 3.2).
\begin{tikzpicture}
  \def\u{0.36}\def\v{0.205}\def\n{12}
  \pgfmathtruncatemacro{\T}{2*\n-2}
  \foreach \t in {0,...,\T} {%
    \foreach \i in {1,...,\n} {%
      \pgfmathtruncatemacro{\s}{\t+\i}
      \pgfmathtruncatemacro{\lc}{\i-1}
      \pgfmathtruncatemacro{\kap}{\n-\i}
      \def\cc{white}%
      \ifnum\t<\lc\relax\else\def\cc{rgAgree}\fi
      \ifnum\s>20
        \ifnum\s<23
          \def\cc{rgInput}%
          \ifnum\kap<7
            \ifnum\s=22 \def\cc{rgCone}\fi
            \ifnum\s=21 \ifnum\kap>0 \def\cc{rgCone}\fi\fi
          \fi
        \fi
      \fi
      \ifnum\s=23 \def\cc{rgReturn}\fi
      \ifnum\s>23 \def\cc{rgRefl}\fi
      \ifnum\t=\T \def\cc{stF}\fi
      \filldraw[fill=\cc,draw=grid!45,line width=0.15pt]
        (\i*\u-\u,-\t*\v) rectangle (\i*\u,-\t*\v-\v);
    }}
  \draw[ink,line width=0.5pt] (0,0) rectangle (\n*\u,-\T*\v-\v);
  \draw[ink,line width=1.3pt] (0,0.06) -- (0,-\T*\v-\v-0.06);
  \draw[ink,line width=1.3pt] (\n*\u,0.06) -- (\n*\u,-\T*\v-\v-0.06);
  % left edge of the dependency cone of the right end at the firing time
  \draw[ink,densely dashed,line width=0.8pt]
    (\n*\u-0.5*\u,-\T*\v-0.5*\v) -- (5.5*\u,-16.5*\v);
  \node[star,star points=5,fill=white,draw=ink,line width=0.4pt,inner sep=0pt,
        minimum size=0.24cm] at (\n*\u-0.5*\u,-\T*\v-0.5*\v) {};
  % relative coordinates at the right end
  \draw[-{Stealth[length=1.6mm]},ink,line width=0.6pt]
    (\n*\u+0.14,-11*\v) -- node[right,font=\scriptsize] {$\tau$} (\n*\u+0.14,-17*\v);
  \draw[-{Stealth[length=1.6mm]},ink,line width=0.6pt]
    (\n*\u,-\T*\v-\v-0.14) -- node[below,font=\scriptsize] {$\kappa$} (\n*\u-5*\u,-\T*\v-\v-0.14);
  \node[font=\scriptsize,anchor=east] at (-0.06,-0.5*\v) {$t=0$};
  \node[font=\scriptsize,anchor=east] at (-0.06,-11.5*\v) {$t=n-1$};
  \node[font=\scriptsize,anchor=east] at (-0.06,-22.5*\v) {$t=2n-2$};
  \node[font=\scriptsize,anchor=south] at (0.5*\u,0.04) {$1$};
  \node[font=\scriptsize,anchor=south] at (\n*\u-0.5*\u,0.04) {$n$};
  % legend
  \begin{scope}[font=\scriptsize]
    \foreach \c/\l [count=\k from 0] in {
        rgAgree/{agrees with the half-line (\cref{lem:agree})},
        rgInput/{input anti-diagonals $t+i=2n-3,\,2n-2$},
        rgCone/{input cells of the cone of $\bigstar$ (\cref{lem:cone})},
        rgReturn/{return chain $t+i=2n-1$ (\cref{lem:return})},
        rgRefl/{rest of the reflected triangle $R_n$},
        stF/{firing row $t=2n-2$}}
      \node[anchor=west,inner sep=0pt] at (5.3,-0.25-0.5*\k) {\swatch{\c}\ \l};
    \node[anchor=west,inner sep=0pt,align=left] at (5.3,-3.35)
      {$\bigstar$\ \ the right end at the firing time; dashed: the left\\
       \phantom{$\bigstar$\ \ }edge of its dependency cone};
  \end{scope}
\end{tikzpicture}

\caption{Anatomy of the line of length $n=12$, coloured by region rather than
by state. The upper part agrees with the half-line (\cref{lem:agree}). The
reflected triangle $R_n$ (\cref{lem:determinacy}) is computed from the two
input anti-diagonals and the borders; its first anti-diagonal is the return
chain (\cref{lem:return}). The right end at the firing time ($\bigstar$)
depends only on the dark input cells (\cref{lem:cone}); the dashed line is the
left edge of its dependency cone. The arrows show the relative coordinates
$\tau$ and $\kappa$ of \cref{sec:right}.}
\label{fig:anatomy}
\end{figure}

\begin{lemma}[determinacy]\label{lem:determinacy}
There is a partial map $\Psi$, depending on $\delta$ only, such that for all
$n\ge2$ and $(\tau,\kappa)\in R_n$ the values $R_n(\tau,\kappa)$ and
$\Psi(J_n,n;\tau,\kappa)$ are both undefined or both defined and equal.
Moreover, $\Psi(J,m;\tau,\kappa)$ depends only on the input cells and the
borders in the backward dependency cone of $(\tau,\kappa)$.
\end{lemma}

In words: the reflected triangle is computed row by row by the local rule,
starting from the input word and the two borders, and each of its cells depends
only on the input cells from which information can reach it.

\begin{proof}
Define $\Psi(J,m;\tau,\kappa)$ for $0\le\kappa\le\tau\le m-1$ by induction on
$\tau$ as $\delta(x_{\kappa+1},x_\kappa,x_{\kappa-1})$, where each $x_{\kappa'}$
is given by the first applicable case: $\ast$ if $\kappa'\in\{-1,m\}$; the lower
or the upper component of $J(\kappa')$ if $\tau-1=\kappa'-2$ or
$\tau-1=\kappa'-1$, respectively; and $\Psi(J,m;\tau-1,\kappa')$ otherwise (then
$0\le\kappa'\le\tau-1$). The value $\Psi$ is undefined if one of the arguments
is $\Fq$ or undefined. For $m=n$ the three arguments are the left, middle and
right predecessors of the cell $(n-1+\tau,n-\kappa)$, and on the input cells
$C_n$ and $C_\infty$ are both undefined or equal by \cref{lem:agree}, so
induction on $\tau$ gives the first claim. Unfolding the definition gives the
second.
\end{proof}

For $(\tau,\kappa)\in R_n$ let $\mathrm{cone}_n(\tau,\kappa)$ be the set of input
cells that are reached from $(\tau,\kappa)$ by repeatedly passing to
predecessors, stopping at input cells and borders. For the right end at the
firing time, $(\tau,\kappa)=(n-1,0)$, this cone covers about the right half of
the two input anti-diagonals.

\begin{lemma}[cone of the right end]\label{lem:cone}
For $n\ge2$, the cone $\mathrm{cone}_n(n-1,0)$ avoids the left border. It
consists of the upper input cell at $\kappa'=0$, both input cells at
$1\le\kappa'\le\lfloor n/2\rfloor$ and, if $n$ is odd, the lower input cell at
$\kappa'=(n+1)/2$.
\end{lemma}

\begin{proof}
A step to a predecessor decreases $\tau$ by one and changes $\kappa$ by at most
one, and the input cells satisfy $\tau'\in\{\kappa'-1,\kappa'-2\}$. A path from
$(n-1,0)$ that reaches an input cell at $\kappa'$ after $r$ steps therefore has
$\kappa'\le r$ and $n-1-r\in\{\kappa'-1,\kappa'-2\}$, so $2\kappa'\le n+1$, and
$2\kappa'\le n$ for upper cells. Conversely, the upper cell at $\kappa'\le n/2$
is reached by $\kappa'$ steps to the left followed by vertical steps, and the
lower cell at $1\le\kappa'\le(n+1)/2$ by $\kappa'-1$ steps to the left,
vertical steps down to $(\kappa'-1,\kappa'-1)$ and one more step to the left,
through cells of $R_n$; the lower cell at $\kappa'=0$ is a predecessor of no
cell of $R_n$. Every cell of $R_n$ on such a path satisfies $\kappa\le\tau\le
n-1-\kappa$, hence $\kappa\le(n-1)/2$, so no such cell is adjacent to the left
border $\kappa=n$.
\end{proof}

\subsection{The Pumping Lemma}

\begin{lemma}[pumping]\label{lem:pumping}
Let $2\le n<n'$, and let $\delta$ synchronize the lines of lengths $n$ and
$n'$ at times $2n-2$ and $2n'-2$. Then $J_n$ and $J_{n'}$ differ on
$\mathrm{cone}_n(n-1,0)$.
\end{lemma}

\begin{proof}
Suppose that they agree there. The relative cell $(n-1,0)$ lies in $R_{n'}$,
and its cone in $R_{n'}$ consists of the same positions, which are away from
the left border $\kappa=n'$. By \cref{lem:determinacy,lem:cone}, the relative
cell $(n-1,0)$ of $R_{n'}$ has the value of the relative cell $(n-1,0)$ of
$R_n$, namely $C_n(2n-2,n)=\Fq$. Thus $C_{n'}(n'+n-2,n')=\Fq$, although
$n'+n-2<2n'-2$.
\end{proof}

We write $\PUMP(n,n')$ for the conclusion of \cref{lem:pumping}. It is a
condition on the half-line alone: the half-line must keep supplying new
information to the right ends of longer and longer lines. If the depth rows
that carry these cones (defined below) were periodic, the input words of two
lengths that differ by a common period would agree on the cone; the next
theorem turns this observation into a bound on where the periodicity can start.

\subsection{Depth Rows and the Speed-1/3 Barrier}

It is convenient to locate the cells of the half-line by their distance behind
the front, the cell $(t,t+1)$. Let
\[
\Phi_t(j):=C_\infty(t,\,t+1-j)\qquad(0\le j\le t)
\]
be the cell at \emph{depth} $j$ behind the front; the front cell has depth $0$.
The local rule becomes
\begin{equation}\label{eq:comoving}
\Phi_{t+1}(j)=\delta\bigl(\Phi_t(j),\,\Phi_t(j-1),\,\Phi_t(j-2)\bigr)
\qquad(0\le j\le t+1),
\end{equation}
with $\Phi_t(-1)=\Phi_t(-2)=\Lq$ (the quiescent cells ahead of the front) and
$\Phi_t(t+1)=\ast$ (the left border). So information only travels to larger
depths. The \emph{depth row} $j$ is the sequence $t\mapsto\Phi_t(j)$, $t\ge j$:
it consists of the cells $(t,t+1-j)$, so it starts in cell $1$ at time $j$ and
then advances by one cell per step, at distance $j$ behind the front. In these
coordinates,
\[
J_n(\kappa)=\bigl(\Phi_{n-3+\kappa}(2\kappa-2),\ \Phi_{n-2+\kappa}(2\kappa-1)\bigr):
\]
the input cells at distance $\kappa$ from the right end lie on the depth rows
$2\kappa-2$ and $2\kappa-1$.

\begin{theorem}[speed-$1/3$ barrier]\label{thm:third}
Let $\delta$ be a minimal-time solution. Every depth row $t\mapsto\Phi_t(j)$ is
eventually periodic; let $\theta(j)$ be the least $t_0\ge j$ such that it is
periodic on $[t_0,\infty)$. For every $n\ge4$ there is a depth $j\le n$ with
$\theta(j)>n+j/2-5/2$. Hence
$\limsup_{j\to\infty}\theta(j)/j\ge 3/2$.
\end{theorem}

\begin{proof}
Row $0$ satisfies $\Phi_{t+1}(0)=\delta(\Phi_t(0),\Lq,\Lq)$, so it is
eventually periodic. By \eqref{eq:comoving}, row $j$ is a finite-state system
driven by the rows $j-1$ and $j-2$; if these are eventually periodic, so is row
$j$. By \eqref{eq:neverfire} all values lie in $W$. Let $P_j$ be a period of
row $j$, fix $n\ge4$ and let $P$ be a common multiple of $P_0,\dots,P_{n-1}$.

We compare the input words of $R_n$ and $R_{n+P}$ on $\mathrm{cone}_n(n-1,0)$.
The input cell at $\kappa=0$ in this cone is the upper one, and it is quiescent
in both triangles by \cref{lem:lightcone}. The other input cells of the cone
have depth $2\kappa-2$ (lower cells) or $2\kappa-1$ (upper cells), at most
$n-1$ by \cref{lem:cone}. Suppose that each of them lies in the periodic part of
its row, that is, $n-3+\kappa\ge\theta(2\kappa-2)$ for the lower cells and
$n-2+\kappa\ge\theta(2\kappa-1)$ for the upper cells. Replacing $n$ by $n+P$
moves each input cell along its depth row by $P$ steps, so $J_{n+P}=J_n$ on the
cone, which contradicts \cref{lem:pumping}. Hence some input cell of the cone,
of depth $j\le n$, satisfies $\theta(j)>n-3+\kappa$. For a lower cell
$\kappa=(j+2)/2$ and for an upper cell $\kappa=(j+1)/2$, so in both cases
$\theta(j)>n+j/2-5/2$.

Suppose now that $\limsup_{j\to\infty}\theta(j)/j<3/2$. Then there are
$\alpha<3/2$ and $\beta$ with $\theta(j)\le\alpha j+\beta$ for all $j$;
since $\theta(j)\ge j$, we may take $\alpha\ge1/2$. For all $j\le n$,
\[
\theta(j)-\frac{j}{2}\le\Bigl(\alpha-\frac12\Bigr)j+\beta\le\Bigl(\alpha-\frac12\Bigr)n+\beta,
\]
and since $\alpha-1/2<1$, the right-hand side is at most $n-5/2$ for $n$ large.
Then $\theta(j)\le n+j/2-5/2$ for all $j\le n$, which contradicts the first
part.
\end{proof}

Depth row $j$ becomes periodic at the cell $(\theta(j),\theta(j)+1-j)$. If
$\theta(j)\ge(3/2-\varepsilon)j$, then $j\le\theta(j)/(3/2-\varepsilon)$, and the
position of this cell satisfies
\[
\theta(j)+1-j\ \ge\ \frac{1-2\varepsilon}{3-2\varepsilon}\,\theta(j)+1,
\]
where the factor $(1-2\varepsilon)/(3-2\varepsilon)$ tends to $1/3$ as
$\varepsilon\to0$. So, for infinitely many $j$, depth row $j$ becomes periodic
only at a time $t$ at which it is close to the cell $t/3$ or beyond: along
these rows, the transient part of the half-line, which starts at the left
border, reaches roughly the cell $t/3$ at time $t$. This is the sense of the
name ``speed-$1/3$ barrier''.

\begin{remark}[an interpretation]\label{rem:division}
In the classical constructions, the line is divided recursively, and the first
division point is where a signal sent by the general meets the echo of the right
end; a \emph{signal} of speed $v$ is a pattern that moves by $v$ cells per
step. A signal of speed $v$ that leaves the general at time $0$ meets the echo,
which leaves the right end at time $n-1$ at speed one, at the cell
$v(2n-1)/(1+v)$. If, as for the two rules simulated below, this signal is the
boundary between the transient and the periodic parts of the depth rows, then
depth row $j$ meets it at the time $(j-1)/(1-v)$, so $\theta(j)\approx j/(1-v)$.
\Cref{thm:third} then says that $v\ge1/3$ asymptotically, so the first division
point, at the cell $v(2n-1)/(1+v)\ge(2n-1)/4$, is asymptotically never in the
half of the line next to the general. Waksman's and Balzer's solutions divide
the line in the middle, with $v=1/3$, and Mazoyer's solution divides it at two
thirds of its length from the general, with $v=1/2$. In simulations of
Mazoyer's solution up to time $2600$, $\theta(j)\in\{2j+1,2j+2\}$ for $1\le
j<700$; for the five-state rule $\delta_{12}$ of \cref{sec:five}, which
synchronizes the lines of lengths $2,\dots,12$, simulations up to time $1400$
give $\theta(j)=\lceil 3j/2\rceil+2$ for $10\le j\le300$, and $\PUMP(n,n')$ holds
for $4\le n<n'\le650$ (\cref{fig:halflines}). This reading is a heuristic and
not a theorem: in general, the boundary of the periodic part need not be a
signal.
\end{remark}

\begin{figure}[t]
\centering
%% Figure 3: the half-lines of Mazoyer's solution and of delta12 for 0 <= t <= 90 (cells from figures.py).
\begin{tikzpicture}[x=0.74mm,y=0.74mm]
  \foreach \file/\title/\dx in {cells_mazoyer_half90/{Mazoyer's solution (six states)}/0,
                                cells_delta12_half90/{$\delta_{12}$ (five states)}/99} {
    \begin{scope}[shift={(\dx,0)}]
      \input{figures/\file}
      \draw[grid,line width=0.3pt] (0,0) rectangle (91,-91);
      \foreach \sl/\st in {2/densely dashed,3/densely dashed} {
        \draw[white,line width=2.0pt] (0,0) -- (91/\sl,-91);
        \draw[ink,\st,line width=0.85pt] (0,0) -- (91/\sl,-91);
      }
      \node[font=\scriptsize,anchor=north,fill=white,inner sep=1pt] at (91/3,-92) {$i=t/3$};
      \node[font=\scriptsize,anchor=north,fill=white,inner sep=1pt] at (91/2+2,-92) {$i=t/2$};
      \node[font=\small,anchor=south] at (45.5,1) {\title};
      \foreach \y in {0,30,60,90}
        \node[font=\scriptsize,anchor=east] at (-0.8,-\y-0.5) {$\y$};
    \end{scope}
  }
  \node[font=\scriptsize,anchor=east] at (-6,-45.5) {$t$};
\end{tikzpicture}

\caption{The half-lines of Mazoyer's six-state solution (left) and of the
five-state rule $\delta_{12}$ (right) for $0\le t\le90$. The dashed lines are
$i=t/2$ and $i=t/3$. The depth rows, that is, the cells at a fixed distance
behind the front, become periodic to the right of the line $i=t/2$ for
Mazoyer's solution and of the line $i=t/3$ for $\delta_{12}$
(\cref{rem:division}); by \cref{thm:third}, no minimal-time solution can do
this to the left of the line $i=t/3$, asymptotically. To the
left of the line $i=t/3$, the half-line of $\delta_{12}$ alternates between
$\Aq$ and $\Lq$, periodically in time and space. \Cref{thm:left} excludes such a
wedge at the left border, and \cref{cor:germ12} makes this explicit for the
line of length $518$.}
\label{fig:halflines}
\end{figure}

\subsection{The Right Band}

The pigeonhole argument behind the pumping lemma also works along the right
border, in the direction of time. The $d$ cells next to the right border form a
finite automaton driven by the column $\kappa=d$, and this automaton has to
reach its all-firing state exactly at the firing time.

\begin{lemma}[right band]\label{lem:rightband}
Let $\delta$ have $k$ states and synchronize the line of length $n$ at time
$2n-2$. Let $1\le d\le\tau_0\le n-2$ and $\ell\ge1$, and suppose that the
column $\tau\mapsto R_n(\tau,d)$ is $\ell$-periodic on $[\tau_0,n-2]$, that is,
$R_n(\tau+\ell,d)=R_n(\tau,d)$ whenever $\tau,\tau+\ell\in[\tau_0,n-2]$. Then
$n-1-\tau_0\le(k-1)^d\ell$.
\end{lemma}

\begin{proof}
For $d-1\le\tau\le n-1$ put $\mathbf b(\tau):=(R_n(\tau,0),\dots,R_n(\tau,d-1))$.
Since $R_n(\tau+1,\kappa)=\delta(R_n(\tau,\kappa+1),R_n(\tau,\kappa),
R_n(\tau,\kappa-1))$ with $R_n(\tau,-1)=\ast$, there is a map $\Gamma$,
depending on $\delta$ and $d$ only, with $\mathbf b(\tau+1)=\Gamma(\mathbf
b(\tau),R_n(\tau,d))$ for $d\le\tau\le n-2$; its arguments lie in $W^d\times W$
because no cell fires before time $2n-2$. At $\tau=n-1$ the band is $\Fq^d$.
Suppose that $n-1-\tau_0>(k-1)^d\ell$. The pairs $(\mathbf
b(\tau),\tau\bmod\ell)$ with $\tau_0\le\tau\le n-2$ take at most $(k-1)^d\ell$
values, so there are $\tau_0\le\tau_1<\tau_2\le n-2$ with $\mathbf
b(\tau_1)=\mathbf b(\tau_2)$ and $\ell\mid h:=\tau_2-\tau_1$. For $0\le s\le
n-2-\tau_2$ the driving cells $R_n(\tau_1+s,d)$ and $R_n(\tau_2+s,d)$ lie in
$[\tau_0,n-2]$ and their times differ by a multiple of the period, so they are
equal, and induction on $s$ gives $\mathbf b(\tau_1+s)=\mathbf b(\tau_2+s)$ for
$0\le s\le n-1-\tau_2$. For $s=n-1-\tau_2$ this reads $\mathbf
b(n-1-h)=\mathbf b(n-1)=\Fq^d$, although $n-1-h<n-1$, a contradiction.
\end{proof}
